\section{Final years, death, burial, and received memory}\label{sec:aug-death}

\subsection{Age, succession, and an unfinished controversy}

By the middle 420s Augustine was an old bishop in ancient terms, still
preaching and dictating while arranging succession. In 426 the people of Hippo
received the presbyter Heraclius as the person who would bear increasing
pastoral responsibility. The minutes preserved as Letter 213 show Augustine
managing the transition publicly rather than treating the see as private
property. He asked that his remaining time be protected for work on Scripture,
but correspondence and controversy continued.

Late in Augustine's episcopate, the \emph{Retractationes} reviewed his books
from the earliest post-conversion works onward; this revision does not assign
its beginning an independently checked exact year. Augustine intended also to review letters and
sermons, a project he did not complete. At his death he was still answering
Julian of Eclanum; the \emph{Unfinished Work against Julian} bears its
condition in its title. The unfinished page is biographically fitting. The
author whom later ages sometimes make into a closed theological monument died
amid argument, correction, and invasion.

The western imperial order in Africa was collapsing. Vandal forces crossed
from Spain in 429 under Geiseric and advanced through the provinces. Possidius
describes devastated churches, displaced populations, and bishops driven from
their sees (\emph{Life} 28). His language is lament and should not be treated
as a neutral military survey, but the catastrophe was real. Hippo became one
of the principal besieged cities in 430. Refugee bishops, including Possidius,
were present with Augustine.

\subsection{Illness and death during the siege}

Possidius says that Augustine fell ill with a fever during the third month of
the siege. In his last days he had penitential psalms copied and placed where
he could see and read them from bed; he asked for privacy except at medical or
meal times (\emph{Life} 31). This scene is a near-contemporary friend's
hagiographic account, but its spirituality agrees with the lifelong
\emph{Confessions}: the bishop dies not as someone beyond repentance but
addressing God in need.

Possidius reports that Augustine lived seventy-six years and spent almost
forty as priest or bishop, then was buried after the sacrifice offered for him.
Ancient inclusive age-counting and Augustine's November birthday mean modern
usage calls him seventy-five. Possidius does \emph{not} state an exact civil
day in the surviving narrative. Prosper of Aquitaine's near-contemporary
\emph{Chronicle}, however, says at the 430 entry that Augustine died
\emph{V kalendas Septembris}---28 August---while answering Julian's books amid
the Vandal assault (MGH AA 9, p.~473). The date is therefore more than a
feast-derived inference, while the compressed chronicle notice still differs
from Possidius's fuller eyewitness-shaped narrative.

Possidius also recounts reports of healings connected with Augustine's prayer,
including a sick man said to have been healed when Augustine laid hands on him
at the request of a visitor (\emph{Life} 29). He presents these as signs of
holiness and notes Augustine's reluctance to claim extraordinary power. This
study reports the witness and its genre; it neither declares the medical event
supernaturally authenticated nor removes it from the life because modern
methods cannot repeat the inquiry.

Hippo did not fall and burn on the day Augustine died. The siege continued,
and the city's later fate belongs to the longer Vandal conquest. Possidius
emphasizes that Augustine's church, library, and books survived the immediate
devastation (\emph{Life} 31). The survival of texts was not automatic. Copying
and movement through North African, Italian, Gallic, and later monastic
networks made the bishop available to ages that no longer inhabited Roman
Hippo.

\subsection{Burial and the limits of the relic itinerary}

Possidius gives burial in the setting of besieged Hippo. Paul the Deacon,
writing in the late eighth century, reports that King Liutprand obtained in
Sardinia remains said to have been moved there earlier because of barbarian
devastation and translated them to Ticinum, or Pavia
(\emph{History of the Lombards} 6.48). This establishes an early-medieval
translation tradition, not the date or witness for the first move, an unbroken
custody chain, or the identity and present condition of particular remains.

Augustine's memory was never limited to one tomb. The 28 August
commemoration, circulation of Possidius's life, and use of Augustine's works
formed a liturgical and textual memory. Later communities claimed him as a
father in different senses, but their institutional histories were not
collated here and must not be backdated into a foundation charter Augustine
never issued.

\subsection{Doctor, contested ancestor, and global reader}

In the Latin West, Augustine became an authority on grace, sacraments,
Trinity, Church, time, memory, interpretation, and political pilgrimage.
Medieval theologians did not read one uniform Augustine; neither did
Reformation and Catholic controversialists. Competing parties cited him on
grace and freedom. Modern philosophers have read his account of time and
interiority; psychologists and memoirists have claimed the
\emph{Confessions}; political thinkers have recruited the two cities. Any
reception that treats him only as a European schoolman, detached from Roman
Africa's soil, conflicts, and multilingual Christianity, suppresses the
setting established by the ancient sources.

Ecclesial reception calls him saint and Doctor without requiring an
unhistorical portrait. His holiness is remembered in conversion, common life,
pastoral endurance, love of Scripture, care for unity and the poor, penitence,
and a mind given to seeking understanding. The same communion can acknowledge
that his defense of coercion, his rhetoric against opponents, and some social
assumptions require criticism. Ecclesial veneration is not a claim that power
left no scar.

\subsection{Visual memory and patronage}

No securely identified contemporary portrait in the checked source set
supplies Augustine's facial appearance. Later art gives him episcopal dress,
books, hearts, and sometimes the child-and-sea motif. This revision has not
established the earliest controlled use of those attributes, so it treats them
only as later visual reception, not dated evidence for his appearance or
possessions. Specific patronage lists are likewise omitted until their
item-level histories and official status can be checked. The book remains a
responsible general attribute because writing, preaching, revision, and the
preservation of a library are directly attested in the life.
